\documentclass[../ICIT-Paper-2024.tex]{subfiles}
\usepackage{hyperref} 
\usepackage{tablefootnote}
\begin{document}
\section{INTRODUCTION}\label{section1}

Decentralized Finance (DeFi)~\cite{werner2022sok} encompasses an ecosystem of decentralized applications (DApps) built on blockchains. One of the most prevalent decentralized applications today is lending protocols or lending pools. According to DefiLlama\footnote{\href{https://defillama.com/}{https://defillama.com/}}, as of May 2024, over 36 million dollars have been locked in lending protocols ranking second in the DeFi ecosystem. Lending pools allow users to deposit their tokens \cite{harvey2021defi} into a liquidity pool~\cite{sun2022liquidity}. These tokens are then utilized to provide loans to other users. The loans are subject to a certain interest rate, depending on the token. The interest accrued from these loans is used to (1) repay deposit interest to depositors and (2) provide a reserve for the lending protocol system \cite{Aave}.

%Tài chính phi tập trung (DeFi) bao gồm một hệ sinh thái các ứng dụng phi tập trung (Dapp) được xây dựng trên các mạng blockchain. Một trong những ứng dụng phi tập trung phổ biến hiện nay là lending protocols hay lending pools. Tính đến tháng 5 năm 2024, có hơn 36 triệu đô được lock trong các lending protocols (theo defilliama) đứng thứ hai trong hệ sinh thái defi. Các lending pools cho phép người dùng gửi các tokens của họ vào nhóm thanh khoản. Lượng tokens này được sử dụng để cung cấp các khoản vay cho người dùng khác. Các khoản vay này sẽ phải chịu một tỷ lệ lãi suất nhất định tùy theo từng loại tokens. Lượng lãi suất cho vay này sẽ được dùng để (1) trả lãi suất gửi cho người gửi tokens và (2) cung cấp một khoản dự phòng cho hệ thống lending protocols. 

Interest rates for loans in lending protocols are significantly higher than those in traditional financial banks. As of May 2024, the lending interest rate for USDC\footnote{\href{https://coinmarketcap.com/currencies/usd-coin/}{https://coinmarketcap.com/currencies/usd-coin/}} on AAVE \footnote{\href{https://app.aave.com/markets/}{https://app.aave.com/markets}} is over 9\%, while Bank of America\footnote{\href{https://www.bankofamerica.com/}{https://www.bankofamerica.com/}} has an average interest rate of over 6\%. Such high interest rates are due to supply not meeting demand, resulting in limited liquidity. The liquidity is provided by depositors and cannot be provided by third parties such as banks or the government, as is possible in traditional finance.

%Lãi suất của các khoản vay trên các lending protocols cao hơn rất nhiều so với ngân hàng trong tài chính truyền thống. Tính đến tháng 5 năm 2024 lãi suất vay usdc của AAVE là hơn 9% trong khi Bank of America thì lãi suất trung bình hơn 6%. Lãi suất cao như vậy là bởi cung không đủ cầu khi lượng tiền cung cấp thanh khoản ít. Lượng thanh khoản cung cấp bởi chính người dùng gửi tiền và không thể  tăng thêm bởi bên thứ ba như ngân hàng hay nhà nước trong tài chính truyền thống

All wallets have the same high interest rates. This leads to two major issues. 
(1) Borrowers are at risk of defaulting and having their collateral liquidated. This occurs when the value of the collateral decreases abruptly. This risk is exacerbated by the ongoing fluctuation of token prices that can drop suddenly at any time. 
For instance, in June 2022, the Ethereum token abruptly dropped in price from \$1800 to \$1000 (according to coingecko\footnote{\href{https://www.coingecko.com/en/coins/ethereum}{https://www.coingecko.com/en/coins/ethereum}}). (2) Lending protocols are unable to optimize profit for the users. This is because all members have different financial strength but having the same rate of borrowed money, and interest rates.

%Lãi suất cao và áp dụng như nhau với mọi ví dẫn đến hai vấn đề lớn. (1) Người vay dễ bị vỡ nợ và bị thanh lý tài sản thế chấp. Điều này xảy ra khi khoản thế chấp bị giảm giá trị đột ngột. Rủi ro này gây ra bởi giá token biến động liên tục và có thể giảm đột ngột bất kỳ lúc nào. (2) Lending protocols không tối ưu được lợi nhuận cho cộng đồng người dùng. Điều này gây ra bởi lượng tiền vay, lãi suất vay, và lãi suất gửi được áp dụng như nhau với toàn bộ cộng đồng. Trong khi những người dùng trong cộng đồng lại có tiềm lực tài chính khác nhau.


Due to these issues, there is an increasing demand for assessing the quality of users' digital wallets. This assessment will determine the user repayment capabilities and classify them accordingly. This process enables lending pools to more effectively establish loan thresholds and interest rates for each participant. Verified and evaluated loans will have a lower risk of liquidation compared to conventional loans presently available.
%Từ hai vấn đề trên, nhu cầu đánh giá chất lượng ví điện tử của người dùng ngày càng cao. Việc đánh giá chất lượng ví điện tử sẽ xác định được khả năng trả nợ của người dùng và phân cấp người dùng. Điều này giúp nhà đầu tư dễ dàng trong việc xác định ngưỡng các khoản vay cũng như lãi suất áp dụng với từng người tham gia lending protocols. Các khoản vay, đã được kiểm định và đánh giá, sẽ có rủi ro bị thanh lý thấp hơn so với các khoản vay thông thường hiện nay.

In traditional finance (TradFi), banks share data and use the FICO Credit Score formula to assess the quality of an individual or organization account. FICO Scores typically range from 300 to 850, with higher scores indicating lower credit risk and vice versa. FICO Scores can be used to evaluate the quality of a digital wallet on DeFi. This is because of the similar features between digital wallets and bank accounts, such as (1) serving as repositories for users' assets, (2) enabling monetary transfers between accounts, and (3) facilitating deposits, collateralization, and asset borrowing.

%Trong giới tài chính truyền thống (TraFi), các ngân hàng chia sẻ dữ liệu cho nhau và sử dụng công thức tính điểm tín dụng (FICO Credit Score) để đánh giá chất lượng tài khoản của một cá nhân hay tổ chức. FICO Scores typically range from 300 to 850 (bảng …), with higher scores indicating lower credit risk and lower scores indicating higher credit risk. FICO Scores có thể được sử dụng để đánh giá chất lượng của một ví điện tử trên Defi. This is because of những đặc điểm tương đồng của ví điện tử và tài khoản ngân hàng như (1) Là nơi lưu giữ tài sản của người dùng, (2) Có thể trao đổi chuyển tiền cho nhau, và (3) Tham gia gửi tiền, thế chấp và vay tài sản.

Over the past few decades, models for optimizing FICO credit score parameters have been developed and presented at scientific conferences. These models often employ regression techniques and deep learning based on information about an individual's outstanding debt, payment history, length of credit usage, mix of credit used, and applications for new credit, etc~\cite{homonoff2021does}. They have been compared and evaluated based on user data from TradFi~\cite{munkhdalai2019empirical}, although this data is difficult to access. This is because it is sensitive and banks only share it privately.
%Trong vài thập kỷ qua, các mô hình tối ưu tham số cho điểm tín dụng FICO Score đã được thực hiện và công bố trong các hội nghị khoa học. Những mô hình này thường thực hiện các kỹ thuật hồi quy và học sâu dựa trên information about an individual's outstanding debt, payment history, length of credit usage, mix of credit used, and applications for new credit, etc. Các mô hình này đã được so sánh và đánh giá dựa trên dữ liệu người dùng trên TraFi dù dữ liệu này rất khó tiếp cận. Điều này là bởi nó là dữ liệu bảo mật và các ngân hàng chỉ chia sẻ private cho nhau. 
This paper proposes a credit score formula within the blockchain domain, utilizing parameters such as total current assets, average total assets, frequency of DApp transactions, number of interacted DApps, types of interacted DApps, reputation of interacted DApps, transaction amount, transaction frequency, account age, number of liquidations, total value of liquidations, loan-to-balance ratio, and loan-to-investment ratio. Our contributions are as follows: (1) the collection, processing, and construction of a dataset comprising 14 characteristics of crypto wallets, and (2) the development of four credit evaluation models based on the FICO score using this dataset, followed by the assessment and comparison of these models.

The remainder of the paper is structured as follows: (\ref{section1}) introduction, (\ref{section2}) related works, (\ref{section3}) methodology, (\ref{section4}) experiment, and (\ref{section5}) conclusion.

\end{document}